\section{The universal numerical interpreter}\label{app:machine}
Here is a self-contained choice of the indexing and interpreter used in Theorem~\ref{output}. It fixes the target independently of the selected reduction. An object program is a finite list of instructions on numbered natural-valued registers. An increment of \(j\) jumps to label \(l\). A conditional decrement of \(j\) jumps to \(l_0\) on zero and to \(l_1\) after a positive decrement. Halt returns register zero. Initially only register zero can be nonzero, with value \(x\). An unoccupied instruction address or a malformed opcode enters an infinite loop. An unconditional jump in an object program uses a permanently zero spare register.

Represent nil by zero and a list cell with head \(a\) and tail \(b\) by
\[
       [a\mid b]=2^a(2b+1).
\]
Every positive integer splits uniquely, and its tail is smaller, so lists are finite. Instruction zero is halt. The increment instruction has code \([0\mid[j\mid l]]\). The conditional decrement has code \([1\mid[j\mid[l_0\mid l_1]]]\). The outer program list of instruction codes has natural code \(e\), with label zero at its head. Require the inner cells in a nonhalting instruction to be positive. The unit program consisting of halt has code \(1\). The program that increments register zero, jumps to label one, and halts has code \([7\mid 1]=384\).

\paragraph{\textit{Interpreter.}}
Reserve two untouched registers \(E,X\) and initialize separate working registers to their values. All other primitive registers start at zero. Set the virtual register list to \([x\mid0]\) and the virtual program counter to zero. The following fixed finite control can be read independently of the accompanying executable.
\begin{enumerate}
\item Copy the program list to a scan register and drop as many cells as specified by the virtual program counter, rejecting a missing instruction by entering a self-jump. Split the selected cell to obtain its instruction head.
\item For halt, take the head of the virtual register list into the designated primitive output register and enter the absorbing halt label. For another instruction, split its fields, check its tag and required inner cells, and retain \(j\) and the branch target or targets.
\item Traverse the virtual register list to index \(j\), saving skipped heads on a reversed temporary list. For register access only, splitting nil gives head and tail zero. This implements implicit trailing zero registers. Split off the addressed value, preserving its tail.
\item Execute the specified operation on this value and set the next virtual program counter. Cons the new value onto the retained tail. Restore saved heads by popping the reversed temporary list and consing them in turn. Return to instruction fetch.
\end{enumerate}

These operations admit a finite primitive expansion. Clear a register by guarded decrements. Transfer its units with a decrement/increment loop. To copy without losing the source, increment the destination and a zero temporary for each consumed unit, then restore the source from that temporary. Double a value with two increments per unit, using an empty temporary. To cons, start with twice the tail plus one and double once per unit of a working copy of the head. To split a positive number, remove pairs of units to form its half and record the parity in finite control. At an even division increment the prospective head and repeat on the quotient. At the first odd division the current quotient is the tail. Even positive quotients strictly decrease before another division. Scratch values at macro boundaries are cleared or consumed. The protected registers \(E,X\) are never decremented, even transiently.

Induction on the virtual program steps shows that a fetch sees exactly the virtual machine's program counter and registers. In a finite access, the traversals are bounded by the supplied address or by the finite reversed prefix. All arithmetic subroutines on a finite input terminate. Thus one defined virtual step takes finitely many primitive steps, the error loop diverges, and a virtual halt returns the designated output. The particular expansion in the companion has 29 registers (including temporaries) and 446 labeled primitive instructions. Its initialization uses only the zero and copy wiring of equation~\eqref{maps}. A choice of different hygienic temporaries or an additional jump would change the target's microstep trace. The queue simulation applies at each of its primitive instructions.

\paragraph{\textit{Universality of the indexing.}}
For completeness, a fixed Turing machine can be compiled into a finite object program. Keep the state and scanned tape symbol in finite control and store the strings strictly left and strictly right of the head in two integer stacks, nearest symbol first. With an alphabet encoded by nonzero digits \(1,\ldots,B-1\) for fixed \(B\), let \(b_\ast\) denote the digit for a blank. Push digit \(d\) as \(Bs+d\); pop by quotient and remainder modulo \(B\). On an empty stack return blank and leave it empty. Blank is also an allowed stored digit. For a right move, push the written symbol on the left stack and pop the right stack into the scanned symbol. A left move uses the symmetric operations. A fixed-\(B\) product is repeated addition. Quotient and remainder are obtained by subtracting units while cycling through \(B\) finite control labels and incrementing a quotient on each completed cycle. Temporary registers implement copies and restores as above. This proves one-step agreement with the Turing machine by induction, with only finitely many object instructions per machine description. A unary input tape is initialized by decrementing a copy of \(x\), and a Turing machine computing a numeric function may return its result as a delimited unary block, counted into register zero. This gives all partial computable functions. Conversely the interpreter's primitive instructions are effective. These facts give the stated universal partial function and its undecidable halting domain, in agreement with classical register-machine universality \cite{register}.
